\chapter{Approche statistique}
\section{Modèle de Poisson -- Maher (1982)}

La loi de Poisson repose sur l'hypothèse que la possession du ballon se transforme 
en but avec une probabilité $p$ faible, mais que le nombre d'occasions est élevé. 
Sous ces conditions, le nombre de buts suit asymptotiquement une loi de Poisson.
On notera que la loi binomiale négative constitue une alternative plus flexible, 
capable de modéliser la surdispersion éventuelle des scores, mais nous retenons 
ici la loi de Poisson conformément à Maher (1982)\cite{maher}.\\

\noindent On note $i$ l'équipe à domicile et $j$ l'équipe extérieure. Le couple de 
scores $(x_{ij}, y_{ij})$ est modélisé comme suit : $X_{ij}$ suit une loi de Poisson 
de moyenne $\alpha_i \beta_j$, $Y_{ij}$ suit une loi de Poisson de moyenne 
$\gamma_i \delta_j$, et $X_{ij}$ et $Y_{ij}$ sont supposés indépendants. 
Les paramètres ont l'interprétation suivante : $\alpha_i$ représente la force 
d'attaque de l'équipe $i$ à domicile, $\beta_j$ la faiblesse défensive de l'équipe 
$j$ à l'extérieur, $\gamma_i$ la faiblesse défensive de $i$ à domicile, et $\delta_j$ 
la force d'attaque de $j$ à l'extérieur.\\

\noindent Pour assurer l'identifiabilité du modèle, on impose les contraintes :
\[
\sum_i \alpha_i = \sum_i \beta_i \qquad \text{et} \qquad \sum_i \gamma_i = \sum_i \delta_i
\]
car les produits $\alpha_i \beta_j$ et $\gamma_i \delta_j$ sont invariants par 
rééchelonnage simultané des deux facteurs.\\

\noindent La log-vraisemblance pour les buts à domicile s'écrit :
\[
\log L(\underline{\alpha}, \underline{\beta}) = \sum_{i} \sum_{j \neq i} 
\left( -\alpha_i \beta_j + x_{ij} \log(\alpha_i \beta_j) - \log(x_{ij}!) \right)
\]

Les conditions du premier ordre donnent :
\[
\frac{\partial \log L}{\partial \alpha_i} = \sum_{j \neq i} 
\left( -\beta_j + \frac{x_{ij}}{\alpha_i} \right) = 0
\]

ainsi, les estimateurs du maximum de vraisemblance $\hat{\alpha}_i$ et $\hat{\beta}_j$ 
vérifient le système implicite :
\[
\hat{\alpha}_i = \frac{\sum_{j \neq i} x_{ij}}{\sum_{j \neq i} \hat{\beta}_j}
\qquad \text{et} \qquad
\hat{\beta}_j = \frac{\sum_{i \neq j} x_{ij}}{\sum_{i \neq j} \hat{\alpha}_i}
\]

Ce système peut être résolu par Newton-Raphson, ce qui conduit à l'approximation :
\[
\hat{\alpha}_i = \frac{\sum_{j \neq i} x_{ij}}{\sqrt{S_x}}, \qquad
\hat{\beta}_j = \frac{\sum_{i \neq j} x_{ij}}{\sqrt{S_x}}, \qquad
S_x = \sum_i \sum_{j \neq i} x_{ij}
\]
où $S_x$ désigne le nombre total de buts marqués à domicile sur l'ensemble des matchs.
Les paramètres $\hat{\gamma}_i$ et $\hat{\delta}_j$ sont estimés de façon analogue 
à partir des $y_{ij}$.

%%%%%%%%%%%%%%%%%%%%%%%%%%%%%%%%%%%%%%%%%%%%%%%%%%%%%%%%%%%%%
\subsection{Adaptation en rolling strict et Lissage Bayésien}

Dans la formulation originale de Maher (1982), les paramètres sont estimés de 
manière statique sur l'ensemble de la saison. Cette approche introduit un 
\textit{data leakage} majeur, car des informations chronologiquement futures sont 
utilisées pour prédire des matchs passés. De plus, la normalisation historique par 
$\sqrt{S_x}$ entraîne un effondrement mathématique des paramètres au fil du temps, ce 
dénominateur croissant indéfiniment avec le nombre total de matchs joués dans la ligue.

Pour remédier à ces limites, nous adoptons un calcul en \textit{rolling strict} : à 
l'instant du match $t$, les paramètres sont estimés uniquement à partir des données 
des matchs joués strictement avant $t$. Afin de stabiliser ces estimations 
face au problème d'absence ou de trop faible volume d'historique en début de saison, 
un lissage bayésien (Bayesian smoothing) est appliqué. 
Ce lissage combine les performances intrinsèques de l'équipe avec un prior 
global, modélisé par la moyenne glissante de la ligue calculée jusqu'à l'instant $t$. 

De plus, pour éviter que les valeurs de $\alpha, \beta$ et $\gamma, \delta$ 
n'explosent, chaque coefficient est divisé par la moyenne de la ligue en cours 

Pour un match opposant l'équipe $i$ (à domicile) à l'équipe $j$ (à l'extérieur) à 
l'instant $t$, les quatre coefficients de force sont définis par :

\[
\alpha_i^{(t)} = \frac{X_i^{(t)} + K \cdot \bar{\mu}_{\text{ligue}}^{(t)}}
{n_i^{(t)} + K} \cdot \frac{1}{\bar{\mu}_{\text{ligue}}^{(t)}}
\]

\[
\delta_i^{(t)} = \frac{Y_i^{(t)} + K \cdot \bar{\nu}_{\text{ligue}}^{(t)}}
{n_i^{(t)} + K} \cdot \frac{1}{\bar{\nu}_{\text{ligue}}^{(t)}}
\]

\[
\beta_j^{(t)} = \frac{Y_j^{(t)} + K \cdot \bar{\mu}_{\text{ligue}}^{(t)}}
{m_j^{(t)} + K} \cdot \frac{1}{\bar{\mu}_{\text{ligue}}^{(t)}}
\]

\[
\gamma_j^{(t)} = \frac{X_j^{(t)} + K \cdot \bar{\nu}_{\text{ligue}}^{(t)}}
{m_j^{(t)} + K} \cdot \frac{1}{\bar{\nu}_{\text{ligue}}^{(t)}}
\]

où :
\begin{itemize}
    \item $X_i^{(t)}$ et $Y_i^{(t)}$ représentent respectivement le cumul des 
    buts marqués et encaissés à domicile par l'équipe $i$ avant le match $t$.
    \item $X_j^{(t)}$ et $Y_j^{(t)}$ représentent respectivement le cumul des 
    buts marqués et encaissés à l'extérieur par l'équipe $j$ avant le match $t$.
    \item $n_i^{(t)}$ et $m_j^{(t)}$ désignent le nombre de matchs joués avant
     $t$, respectivement à domicile par $i$ et à l'extérieur par $j$.
    \item $\bar{\mu}_{\text{ligue}}^{(t)}$ et $\bar{\nu}_{\text{ligue}}^{(t)}$ sont 
    les moyennes mobiles des buts marqués par match dans l'ensemble de la ligue 
    (respectivement à domicile et à l'extérieur) calculées sur l'historique strictement 
    avant $t$.
    \item $K$ est l'hyperparamètre de lissage, fixé ici à $K = 5$. Il fait office 
    de force de rappel en simulant l'équivalent de 5 matchs fictifs aux performances 
    moyennes, ce qui s'accorde de manière cohérente avec la fenêtre glissante 
    utilisée lors du prétraitement des données.
\end{itemize}

Les espérances de buts finales ($\mu^{(t)}$ pour l'équipe à domicile et $\nu^{(t)}$ 
pour l'équipe à l'extérieur) utilisées pour alimenter les lois de Poisson 
indépendantes sont alors réajustées par l'avantage moyen du terrain et s'écrivent :
\[
\mu^{(t)} = \alpha_i^{(t)} \cdot \beta_j^{(t)} \cdot \bar{\mu}_{\text{ligue}}^{(t)} 
\qquad \text{et} \qquad \nu^{(t)} = \gamma_j^{(t)} \cdot \delta_i^{(t)} \cdot 
\bar{\nu}_{\text{ligue}}^{(t)}
\]
Cette formulation garantit que $\mu^{(t)}$ et $\nu^{(t)}$ conservent rigoureusement 
leur sens statistique d'espérances de buts par match à l'instant $t$, tout en 
éliminant les anomalies sur les scores extrêmes.

\subsection{Métriques d'évaluation des buts}
Erreur Quadratique Moyenne : $MSE = \frac{1}{n} \sum(y - \hat{y})^2$


L'Erreur Quadratique Moyenne (MSE) mesure la moyenne des carrés des écarts entre les 
valeurs prédites et les valeurs réelles. 
En élevant au carré les erreurs de prédiction, la MSE pénalise les erreurs importantes. 
Cette propriété la rend particulièrement 
utile dans des contextes où les grandes erreurs doivent être évitées à tout prix.\\


Erreur Absolue Moyenne (MAE) :
$MAE= \frac{1}{n} \sum|y - \hat{y}|$ 

L'Erreur Absolue Moyenne (MAE) est une métrique qui mesure, en moyenne, la 
magnitude des erreurs entre les valeurs prédites par le modèle et les valeurs 
réelles observées, sans tenir compte de la direction de l'erreur. 
La MAE est appréciée pour sa simplicité et son interprétation intuitive. Elle donne 
une estimation directe de l'erreur moyenne sur l'ensemble des prédictions. Étant 
donné qu'elle utilise la valeur absolue des erreurs, elle est robuste face aux 
valeurs aberrantes, contrairement à d'autres métriques qui amplifient les erreurs.\\

mean\_poisson\_deviance\\
Pour une observation :
\[
D(y, \hat{y}) = 2 \left( y \log\left(\frac{y}{\hat{y}}\right) - (y - \hat{y}) \right)
\]
Pour $n$ observations :
\[
\text{mean\_poisson\_deviance} = \frac{1}{n} \sum_{i=1}^{n} 2 \left( y_i \log\left(\frac{y_i}{\hat{y}_i}\right) - (y_i - \hat{y}_i) \right)
\]

Cette formule mesure l'erreur entre les valeurs réelles $y$ et les prédictions 
$\hat{y}$ pour des données de comptage (0,1,2,\dots). Plus la valeur est petite, 
meilleur est le modèle.

%%%%%%%%%%%%%%%%%%%%%%%%%%%%%
%%%%%%%%%%%%%%%%%%%%%%%%%%%%%
\section{Approche statistique de Dixon et Coles}
Dans la continuité du modèle de Maher, Dixon et Coles ont appliqué une fonction $\tau$ pour gérer 
les petits scores $x$-$y$ avec $x,y \in \left\{ 0,1 \right\} $.
\[
\tau(x, y) = 
\begin{cases} 
    1 - \mu \nu \rho & \text{si } x=0, y=0 \\ 
    1 + \mu \rho     & \text{si } x=0, y=1 \\ 
    1 + \nu \rho         & \text{si } x=1, y=0 \\ 
    1 - \rho             & \text{si } x=1, y=1 \\ 
    1                    & \text{sinon}
\end{cases}
\]

\noindent Nous avons ainsi, $$L(\alpha, \beta, \gamma, \delta, \rho) = \tau(x, y, \mu, \nu, \rho) \times 
\frac{\mu^x e^{-\mu}}{x!} \times \frac{\nu^y e^{-\nu}}{y!}$$
D'où : $\log L=\log(\tau)+x\log(\mu)-\mu-\log(x!) + y\log(\nu)-\nu-\log(y!)$. 
Nous cherchons $\rho$ qui maximise cette log-vraisemblance. Nous avons donc, 
$$\frac{\partial}{\partial\rho}\log L = \frac{\tau'}{\tau} $$ 
or 
\[\tau'(x,y)=
\begin{cases} 
    - \mu \nu  & \text{si } x=0, y=0 \\ 
    \mu        & \text{si } x=0, y=1 \\ 
    \nu        & \text{si } x=1, y=0 \\ 
    - 1        & \text{si } x=1, y=1 \\ 
    0          & \text{sinon}
\end{cases}
\]

$\rho^*$ vérifie $\partial_\rho \log L(\rho^*) = 0$
Dans notre cas, étant donné $N$ matchs, nous avons : 
$$\partial_\rho \log L = \sum_{1}^{N}\frac{\tau'(x_k,y_k)}{\tau(x_k,y_k)}$$
Autrement dit,
$$ \partial_\rho \log L =
\sum_{k \in \{0-0\}} \frac{-\mu_k \nu_k}{1 - \mu_k \nu_k \rho^*} 
+ \sum_{k \in \{0-1\}} \frac{\mu_k}{1 + \mu_k \rho^*} 
+ \sum_{k \in \{1-0\}} \frac{\nu_k}{1 + \nu_k \rho^*} 
+ \sum_{k \in \{1-1\}} \frac{-1}{1 - \rho^*}$$



\section{Prédiction de l'issue du match}

Comme nous l'avons précisé en introduction de ce rapport, deux prédictions sont majoritaires. L'une basée sur 
le nombre de buts, l'autre sur l'issue du match. Les modèles statistique implémentés se basent sur le nombre de buts. 
En effet, la distribution de Poisson décrit l'évolution du nombre de buts. Cependant, les modèles de machine 
learning utilisés, eux, prédisent l'issue des matchs, nous devons donc effectuer cette transition afin de pouvoir 
comparer l'approche ML avec l'approche statistique. 

Pour cela, selon Maher, les buts marqués par l'équipe A et par l'équipe B sont indépendants. Par suite, la probabilité 
$P(X=x, Y=y) = P(X=x) \times P(Y=y)$ ce qui nous permet de construire la matrice $M$ telle que 
$M_{i,j} = P(X=i, Y=j)$ figure \ref{fig:matricMaher}. Ainsi, la probabilité que l'équipe A gagne correspond à la probabilité qu'elle 
marque plus de buts que l'équipe B. Autrement dit, la victoire de l'équipe à domicile correspond aux coefficients de 
la matrice tel que $i>j$. Le match nul correspond à la diagonale, et la victoire extérieur correspond à la matrice triangulaire 
supérieure $i<j$ figure \ref{fig:matrixMaherSeparation}.

\begin{figure}[H]
    \centering
    % --- Minipage Gauche (Matrice de Maher) ---
    \begin{minipage}[b]{0.48\textwidth}
        \centering
        \includegraphics[width=\textwidth]{../../plot/exemple_matrix.png}
        \caption{Matrice de Maher}
        \label{fig:matricMaher}
    \end{minipage}
    \hfill % Espace élastique pour repousser les deux blocs sur les bords
    % --- Minipage Droite (Séparation en 3 matrices) ---
    \begin{minipage}[b]{0.48\textwidth}
        \centering
        \includegraphics[width=\textwidth]{../../plot/ex_matrix_DC.png}
        \caption{Matrice de Dixon \& Coles}
        \label{fig:matricDC}
    \end{minipage}
\end{figure}

\begin{figure}[H]
    \centering
    % --- Image du Haut ---
    \includegraphics[width=0.85\textwidth]{../../plot/ex_matrix.png}
    \caption{Modèle de Maher : Séparation en 3 matrices distinctes}
    \label{fig:matrixMaherSeparation}
    
    \vspace{0.5cm} % Ajuste l'espace vertical entre les deux images si besoin
    
    % --- Image du Bas ---
    \includegraphics[width=0.85\textwidth]{../../plot/ex_matrix_DC_2.png}
    \caption{Modèle de Dixon-Coles : Séparation en 3 matrices distinctes}
    \label{fig:matrixDCSeparation}
\end{figure}

Dans le cas de Dixon \& Coles, l'hypothèse d'indépendance n'est pas tout a fait la même dans le cas des score faibles. 
Pour traiter ce cas, ils ajoutent un facteur $\tau$, la matrice $M$ devient donc : 
$M_{i,j}=\tau(i,j)\times P(X=i)\times P(Y=j)$.

